\section{The three prices: the resolution grid and finite shots}
\label{sm:prices}

\paragraph{The support exponent and its fit range.} Restricting the fit of Eq.~\eqref{eq:supportexp} to
$L=4$--$12$, the five sizes Table~\ref{tab:fracgrid} covers, gives $1.082$ at $R^{2}=0.998$; we quote
the fit range with the exponent throughout, because it matters at the second decimal. The same caution
applies to the object fitted: $1.082$ is the fit to the published points, while the $1.090$ in the
bottom block of Table~\ref{tab:fracgrid} is the fit to that table's own coarser replication grid at the
same resolution. They are two fits to two datasets and the difference is theirs, not a rounding.

\paragraph{The two exponents are one exponent.} It would be convenient to report that the support
exponent, unlike the fraction exponent, is resolution-independent. It is not, and the arithmetic is
immediate: the fraction \emph{is} the support divided by the sector dimension --- exactly, to the last
digit, at all six sizes --- and the sector dimension does not depend on $\eta$. The two exponents
therefore differ by a constant,
$b_{|\mathcal S|}(\eta)-b_{\rm frac}(\eta)=b_{\dim}=1.291$ on $L=4$--$12$, and they move together:
across the decade $\eta=0.05\to0.50\,t$ both shift by $0.148$ in absolute value
(Table~\ref{tab:fracgrid}). There is one resolution-dependent exponent in this analysis, and quoting
the fraction and the support separately reports it twice.

What differs is the base it is measured against, and that is not cosmetic. On the fraction exponent
$0.148$ is three quarters of the quantity itself: $e^{-0.13\,L}$ becomes $e^{-0.27\,L}$, and the
headline changes qualitatively with a choice the model does not make. On the support exponent the same
$0.148$ carries it from $1.165$ to $1.017$---a seventh of its own value, $13\%$ measured against the
larger end and $15\%$ against the smaller---and ``exponential in $L$ with an exponent near one''
survives at every resolution in the decade. We therefore report the support, and we report it for two reasons, neither of which is
stability: it is the absolute resource a device has to populate, and it is not divided by a sector
dimension that grows for a reason---the combinatorics of the Hilbert space---that has nothing to do
with the method. The fraction is reported as the derived, resolution-dependent quantity that it is,
with its full grid attached.

\begin{table*}[tp]\centering\small
\caption{\textbf{Sector fraction required for relative $L_1<0.05$, versus system size and resolution.}
Same ground state, same time-evolution ranking (protocol $K=18$, $dt=0.5$, $16$ substeps, $19$
snapshots; \texttt{scaling\_lanczos\_mf.py}), same threshold, same depth $n_\ell=150$; only $\eta$
changes. The scan over subspace size in this replication is coarser than the bisection that produced
the published points, and returns the smallest scanned fraction meeting the threshold, so the
$\eta=0.15\,t$ column is an upper bound on them: against the published $0.9167$, $0.8200$, $0.5561$,
$0.3625$, $0.1700$ it reads $0.917$, $0.820$, $0.560$, $0.380$, $0.180$, agreeing at the two smallest
sizes and standing $0.7\%$, $4.8\%$ and $5.9\%$ high at the three largest. The same granularity is why
the fraction exponent of this column, $-0.2012$, is not identical to the $-0.2093$ obtained by fitting
the published points themselves over $L=4$--$12$. Bottom block: least-squares exponents. \emph{Both}
fit windows are given for the support, because the choice moves the third digit and this work
criticizes exactly that elsewhere in this section. The fraction exponent moves by a factor $2.2$ over
the decade, from $-0.1266$ to $-0.2743$, while the support exponent moves from $1.165$ to $1.017$; in
absolute value both move by $0.148$, which is the whole content of the
comparison.}\label{tab:fracgrid}
\begin{tabular}{lcccccccc}
\toprule
 & \multicolumn{8}{c}{$\eta/t$}\\
\cmidrule(lr){2-9}
$L$ & $0.05$ & $0.075$ & $0.10$ & $0.15$ & $0.18$ & $0.25$ & $0.35$ & $0.50$\\
\midrule
$4$  & $0.917$ & $0.917$ & $0.917$ & $0.917$ & $0.917$ & $0.917$ & $0.917$ & $0.875$\\
$6$  & $0.900$ & $0.890$ & $0.870$ & $0.820$ & $0.800$ & $0.730$ & $0.680$ & $0.630$\\
$8$  & $0.740$ & $0.680$ & $0.640$ & $0.560$ & $0.540$ & $0.480$ & $0.420$ & $0.380$\\
$10$ & $0.520$ & $0.480$ & $0.440$ & $0.380$ & $0.340$ & $0.260$ & $0.220$ & $0.200$\\
$12$ & $0.340$ & $0.280$ & $0.240$ & $0.180$ & $0.160$ & $0.140$ & $0.120$ & $0.100$\\
\midrule
fraction exponent, $L{=}4$--$12$ & $-0.1266$ & --- & $-0.1681$ & $-0.2012$ & --- & $-0.2395$ & --- & $-0.2743$\\
$R^{2}$                          & $0.904$   & --- & $0.918$   & $0.929$   & --- & $0.970$   & --- & $0.982$\\
support exponent, $L{=}4$--$12$  & $1.165$   & --- & $1.123$   & $1.090$   & --- & $1.052$   & --- & $1.017$\\
$R^{2}$                          & $0.999$   & --- & $0.999$   & $0.998$   & --- & $0.999$   & --- & $0.999$\\
support exponent, $L{=}6$--$12$  & $1.136$   & --- & $1.088$   & $1.053$   & --- & $1.022$   & --- & $0.992$\\
\bottomrule
\end{tabular}
\end{table*}

The headline these numbers offer---that the fraction needed for a fixed spectral accuracy falls
from $0.92$ at $L=4$ to $0.08$ at $L=14$, as $e^{-0.25\,L}$---is tunable, and no headline of that form
is claimed here. Over the decade of $\eta$ the fraction exponent moves by a factor $2.2$, from
$-0.1266$ at $\eta=0.05\,t$ to $-0.2743$ at $\eta=0.50\,t$. At $\eta=0.50\,t$ the fraction at $L=12$
is already $0.10$, a quarter above the $L=14$ value at the resolution used throughout this work and
reached two sizes earlier: any headline of this form is reachable by choosing the broadening.

Nor is the fraction exponential even at fixed $\eta$. The local logarithmic slopes between consecutive
published points are $-0.056$, $-0.194$, $-0.214$, $-0.379$, $-0.377$---a factor $6.8$ across five
intervals, steepening monotonically over the first four. A single-exponential fit over $L=4$--$14$, with $R^{2}=0.94$, is a straight line drawn through a curve of increasing slope, and the fitted exponent depends on where the fit stops:
$L=4$--$10$ gives $-0.159$, $L=4$--$12$ gives $-0.209$, and the $-0.2477$ of the fit over $L=4$--$14$ requires reaching
$L=14$.

The mechanism that makes the fraction move with $\eta$ is available without derivation, and
Sec.~\ref{sec:moments} states it as Conjecture~\ref{conj:geom} rather than as a theorem:
the polynomial-approximation route to a geometric rate in the Krylov order places the kernel poles at
$\omega\pm i\eta$, so any such rate approaches $1$ as $\eta\to0^{+}$ and any prefactor diverges
(Sec.~\ref{app:moments:conj}). At finer resolution more Krylov order---hence more of the sector---is
needed for the same $L_1$ error. We use that as a mechanism and not as a bound: no $C(\eta)$ or
$\rho(\eta)$ is exhibited anywhere in this work, and none is needed to read Table~\ref{tab:fracgrid},
which is a measurement.

One further constraint has to be declared rather than inherited. The reference against which the
threshold is applied is itself a finite-depth Lanczos computation, recomputed at the same depth as the
subspaces so that the comparison is internally consistent ($n_\ell=150$;
\texttt{scaling\_lanczos\_mf.py}). Two consequences follow, and both belong in the open. First, the
published relative $L_1$ is a distance between two depth-$150$ objects, not between the reconstruction
and the exact $A$. Second, that reference drifts with depth, and the drift is resolution-dependent:
between $n_\ell=100$ and $n_\ell=200$ at $L=12$ the relative $L_1$ change of the reference
\emph{alone} is $11.1\%$ at $\eta=0.05\,t$, $3.0\%$ at $\eta=0.10\,t$, $0.93\%$ at $\eta=0.15\,t$ and
$0.05\%$ at $\eta=0.25\,t$. At the depth used here the protocol is self-consistent only for
$\eta\gtrsim0.15\,t$---which is exactly where the published column sits, at the margin, and it is
\emph{below} that margin, at $\eta=0.10\,t$, that Fig.~\ref{fig:spin} is drawn.

The fraction is therefore reported as a declared window rather than as a number: only for
$\eta\ge0.15\,t$ at $n_\ell=150$, with the full $\mathrm{frac}(L,\eta)$ grid of
Table~\ref{tab:fracgrid} published beside it.

\paragraph{The shot budget, in full.} Populating a subspace by measurement is a coupon-collector problem on the Born distribution of the
time-evolved seed, so the budget exceeds the support it buys. Table~\ref{tab:shots} reports that budget
\emph{measured}, by multinomial sampling of the declared protocol, at the five sizes where the
measurement is affordable; only $L=14$ carries no measurement. Two independent runs contribute, and we
name them as two rather than merging them: run~A draws eight seeds at $L=4$, $6$ and $8$, run~B draws
four seeds at $L=4$, $6$, $8$ and $12$ and six at $L=10$, with different seeds and a different budget
grid. Where they overlap they agree---$2.43$ against $2.44\times10^{4}$ at $L=6$ and $2.70$ against
$2.61\times10^{5}$ at $L=8$, that is to $0.2\%$ and $3.2\%$---which is why the two larger rows are
reported as measurements and not as extrapolations of the smaller ones. Every budget is for \emph{one}
orbital seed and \emph{one} branch; the momentum-resolved observable consumes the full set of orbital
seeds in both $(N{\pm}1)$ sectors, sixteen channels at $L=8$.

\paragraph{The published curve and finite shots, in full.} Precision about the label matters more here than anything else. The ranking that produced the
published fractions is \texttt{argsort} applied to the Born distribution of the time-evolved state
accumulated over the $19$ snapshots of the protocol (\texttt{scaling\_lanczos\_mf.py},
\texttt{stage\_rank})---that is, to the distribution the device would itself measure, in the limit of
infinitely many shots. It is \emph{not} a clairvoyant ordering by exact eigenvector amplitudes; the
repository contains such an oracle elsewhere (\texttt{amortized\_recovery.py}) and labels it as one.
But it is also not a finite-shot reconstruction, and to describe it as \emph{sampled} without a shot
count is to mislabel it. The published curve is the $T\to\infty$ limit of the declared protocol, and it is
labelled that way from here on.

\begin{table*}[tp]\centering\small
\caption{\textbf{Finite shots, with the budget attached.} Multinomial sampling of the declared
protocol, eight independent seeds, $\eta=0.15\,t$; $\pm$ is the sample standard deviation across
seeds, not the standard error of the mean. Top block: the sector fraction the finite-shot
reconstruction needs to reach relative $L_1<0.05$, and the budget at which it reaches it. Bottom
block: at \emph{matched} subspace size, the error of the infinite-shot ordering against the error of
finite shots at the stated budget, and the ratio between them.}\label{tab:finiteshot}
\begin{tabular}{lccc}
\toprule
 & $L=4$ & $L=6$ & $L=8$\\
\midrule
fraction, $T\to\infty$ (published)     & $0.9167$ & $0.8200$ & $0.5561$\\
fraction, finite shots                 & $0.9167\pm0.0000$ & $0.8571\pm0.0135$ & $0.6275\pm0.0031$\\
budget $T$ at which it is reached      & $(7.8\pm2.0)\times10^{2}$ & $(2.43\pm0.34)\times10^{4}$ & $(2.70\pm0.13)\times10^{5}$\\
\midrule
matched $|\mathcal S|$ (size the sampler returned) & $22$ & $247$ & $2331$\\
budget $T$ at that size                & $7.8\times10^{2}$ & $1.74\times10^{4}$ & $2.04\times10^{5}$\\
relative $L_1$, $T\to\infty$ ordering  & $1.4\times10^{-14}$ & $0.0469$ & $0.0362$\\
relative $L_1$, finite shots           & $1.4\times10^{-14}$ & $0.0702\pm0.0115$ & $0.0709\pm0.0042$\\
ratio                                  & $1.00$ & $1.50\pm0.25$ & $1.96\pm0.12$\\
\bottomrule
\end{tabular}
\end{table*}

The published curve is therefore a strict lower bound on the finite-shot cost. At $L=4$ the two agree
because the sampler exhausts the sector; at $L=6$ and $L=8$ the finite-shot fraction exceeds the
published one by $0.037$ and $0.071$. That is a direction across three sizes, one of which is
degenerate, and we state it as a direction and do not fit an exponent to it---which is the same
objection this section raises against the fraction headline, applied to our own finite-shot curve. At the
published subspace itself, $L=8$ and $|\mathcal S|_\eta=2180$, the infinite-shot ordering gives
relative $L_1=0.0490$, just inside the threshold, against $0.1006\pm0.0036$ for finite shots at
$T=1.47\times10^{5}$: a factor $2.05\pm0.07$, and the reason the finite-shot fraction has to rise to
$0.63$.

The finding that runs the other way is the more useful one, and it is what this subsection closes on:
the penalty is budgetary, not structural, and it extinguishes with a budget that is affordable at this
size. On the same pair as the last row of Table~\ref{tab:finiteshot}---$L=8$,
$|\mathcal S|_\eta=2331$---raising the budget from $T=2.04\times10^{5}$ to $T=2.60\times10^{6}$, a
factor $12.7$, takes the ratio of the finite-shot error to the infinite-shot-ordering error from
$1.96\pm0.12$ to $1.06\pm0.01$, through $1.51$, $1.18$ and $1.09$ at the intermediate budgets. No
sampling penalty in this family can be quoted without its budget, because the penalty is a statement
about $T$ and not about the method. We therefore report finite-shot fractions only with $T$ and the
seed count attached, and we report the infinite-shot ordering for what it is---the floor that any
budget approaches from above, at the price Table~\ref{tab:shots} now states.
